% ===== sec/05_chomsky_classification.tex =====
\section{Chomsky Classification of Agents}
\label{sec:chomsky}

We now classify trace languages by Chomsky type and exhibit, for
each type, the canonical agent shape from the empirical
\textsc{Spaceship Titanic} system.  The classification is the
literal one of Section~\ref{sec:preliminaries}; there is no
new mathematics here, only the application of an old hierarchy to
a new object.  The novelty, if any, is the identification of
which sub-agents in the empirical CSVs inhabit which class.

\subsection{Type-3: regular trace languages}

\begin{definition}[Type-3 agent]
\label{def:type3agent}
An agent $\mathcal{A}$ is \emph{Type-3} if $L(\mathcal{A}) \in
\mathcal{L}_3$, equivalently, if there exists a finite
automaton $M$ with $L(M) = L(\mathcal{A})$.
\end{definition}

In the empirical system, the Type-3 agents are
\texttt{tools\_index} and \texttt{unit\_test\_runner}.  Both have
trace languages over a small finite alphabet whose strings are
generated by a finite-state controller:
\begin{itemize}
\item \texttt{tools\_index}: the trace language is the regular
expression
\[
(\textsf{load\_schema})\cdot(\textsf{compose\_tool\_text})\cdot
(\textsf{embed\_tool\_text})\cdot(\textsf{score\_against\_index})
\cdot \dots
\]
read straight off the column.  No back-edge appears.  No stack
discipline is required.
\item \texttt{unit\_test\_runner}: the trace is a sequence of
\textsf{test\_*} keywords drawn from the finite test catalogue.
The deciding finite automaton is the Aho--Corasick DFA over the
catalogue~\cite{aho1975corasick}.
\end{itemize}

The Data-Driven verifier $\mathcal{V}$ of this paper is realised
as the union of these two columns, compiled into a single
Aho--Corasick DFA.  Membership in $L(\mathcal{V})$ is decided in
$O(n)$ time and $O(1)$ space per candidate trace, by Rabin and
Scott's classical
construction~\cite{rabin1959finite,aho1975corasick}.

\subsection{Type-2: context-free trace languages}

\begin{definition}[Type-2 agent]
\label{def:type2agent}
An agent $\mathcal{A}$ is \emph{Type-2} if $L(\mathcal{A}) \in
\mathcal{L}_2 \setminus \mathcal{L}_3$, equivalently, if there
exists a pushdown automaton $M$ with $L(M) = L(\mathcal{A})$ and
$L(\mathcal{A})$ is not regular.
\end{definition}

In the empirical system, the Type-2 sub-agents are
\texttt{planner\_agent} and \texttt{reviewer\_agent}.  Both exhibit nested structure that is naturally modelled with a stack.
The planner emits a step-list whose nesting depth is determined
by the EDA bundle:
\begin{verbatim}
emit_step(id; title; instructions; tools)
  -> emit_step(...) -> ...
emit_success_criteria
emit_plan_json
\end{verbatim}
and the reviewer's score-and-clamp loop pushes a configuration
onto an implicit stack each time a sub-test fails:
\begin{verbatim}
collect_failed_tests -> compose_reviewer_prompt
  -> clamp_score_to_2_on_unit_failure
  -> clamp_score_to_1_5_range
  -> emit_review_result(...).
\end{verbatim}
The stack discipline is implicit but recoverable: the trace
language of \texttt{planner\_agent} \emph{matches Dyck-language
shape} \cite{hopcroft2006automata} when one identifies opening
\textsf{compose\_planner\_prompt} and closing
\textsf{emit\_plan\_json} as bracket symbols.

\subsection{The indexed-grammar zone (between Type-2 and Type-1)}

Aho's indexed grammars~\cite{aho1968indexed} occupy a class
$\mathcal{L}_{\mathrm{indexed}}$ strictly between $\mathcal{L}_2$
and $\mathcal{L}_1$.  Their distinguishing feature is a finite
\emph{stack of indices} attached to each nonterminal during a
derivation.

In the empirical system, the indexed-zone sub-agent is
\texttt{sop\_orchestrator}, whose trace language exhibits a
finite-index repair stack: the tokens
\texttt{cap\_iterations\_at\_max(=3)} and
\texttt{retry\_phase\_on\_low\_score} together form an indexed
loop whose nesting is bounded by the constant 3.  This loop resembles the bounded-index bookkeeping that indexed grammars capture.  We do not prove non-context-freeness here; we use the indexed zone as a modelling abstraction for the coupled iteration and phase counters.

\subsection{Type-1: context-sensitive trace languages}

\begin{definition}[Type-1 agent]
\label{def:type1agent}
An agent $\mathcal{A}$ is \emph{Type-1} if $L(\mathcal{A}) \in
\mathcal{L}_1 \setminus \mathcal{L}_{\mathrm{indexed}}$,
equivalently, if there exists a linear-bounded automaton $M$
with $L(M) = L(\mathcal{A})$ and $L(\mathcal{A})$ is not in the
indexed class.
\end{definition}

The Type-1-style sub-agents are \texttt{task\_orchestrator} and
\texttt{kernel\_pipeline}.  Their trace languages exhibit
linear-bounded workspace growth: each step's emitted string is
bounded by a constant times the input length, and the
\texttt{compute\_dependency\_graph(prev\_stage\_to\_next\_stage)}
operation is, by construction, an LBA-style traversal of the
phase graph that visits each cell at most a constant number of
times.

The \texttt{kernel\_pipeline}'s trace contains
\textsf{generate\_kernel\_packages(4\_copies)},
\textsf{write\_packages\_to\_disk}, and
\textsf{sequential\_push\_on\_vm(stage\_order)} --- a strictly
linear workspace expansion in the number of stages.  No string
in $L(\texttt{kernel\_pipeline})$ exceeds a constant multiple of
the input description's length, which is the LBA hallmark.

\subsection{Type-0: recursively enumerable trace languages}

\begin{definition}[Type-0 agent]
\label{def:type0agent}
An agent $\mathcal{A}$ is \emph{Type-0} if $L(\mathcal{A}) \in
\mathcal{L}_0 \setminus \mathcal{L}_1$, equivalently, if there
exists a Turing machine $M$ with $L(M) = L(\mathcal{A})$ and
$L(\mathcal{A})$ is not context-sensitive.
\end{definition}

The components for which we use an idealised Type-0 envelope are
\texttt{mlebench\_solver}, the \texttt{developer\_agent} (in its
idealised code-generation mode), and the \texttt{kaggle\_trainer}.
The \texttt{kaggle\_trainer} column is not a proof of Type-0
membership; it is an empirical witness that the implementation used
a recursively structured enumeration of training operations
\textsf{bootstrap\_indices}, \textsf{fit\_model\_on\_subset},
\textsf{score\_train\_and\_val},
\textsf{accumulate\_oof\_predictions}, and
\textsf{track\_best\_episode\_by\_val\_score}, parameterised by
loop bounds (\textsf{for\_each\_fold\_for\_each\_episode}) that
are themselves outputs of the agent's own deliberation.

The \emph{Goal-Driven} agent of the architecture is therefore
modelled by an idealised recursive-enumeration envelope under the
orchestration relation $R$ given in Section~\ref{sec:methodology}.
The concrete implementation remains bounded; undecidability enters
only as a property of the idealised envelope, not as an empirical
claim about the finite run.

\subsection{Summary table}

\begin{table}[ht]
\centering
\caption{Language-class abstractions suggested by the empirical sub-agent traces.  Strict class membership is asserted only where an explicit recogniser is supplied; the remaining rows are modelling hypotheses.}
\label{tab:chomsky-empirical}
\resizebox{\columnwidth}{!}{%
\begin{tabular}{l|l|l}
\hline
\textbf{Sub-agent} & \textbf{Abstraction} & \textbf{Evidence} \\ \hline
\texttt{tools\_index}        & Type-3 & FA / DFA \\
\texttt{unit\_test\_runner}  & Type-3 & FA / Aho--Corasick DFA \\
\texttt{reader\_agent}       & Type-3 & FA \\
\texttt{summarizer\_agent}   & Type-3 & FA \\ \hline
\texttt{planner\_agent}      & Type-2? & PDA hypothesis \\
\texttt{reviewer\_agent}     & Type-2? & PDA hypothesis \\ \hline
\texttt{sop\_orchestrator}   & Indexed? & Indexed-grammar hypothesis
                                          (Aho \cite{aho1968indexed}) \\ \hline
\texttt{task\_orchestrator}  & Type-1? & LBA hypothesis \\
\texttt{kernel\_pipeline}    & Type-1? & LBA hypothesis \\ \hline
\texttt{research\_agent}     & Type-0? & TM hypothesis \\
\texttt{developer\_agent}    & Type-0? & TM hypothesis \\
\texttt{mlebench\_solver}    & Type-0? & TM hypothesis \\
\texttt{kaggle\_trainer}     & Type-0? & TM hypothesis (emergent) \\ \hline
\end{tabular}%
}
\end{table}

\subsection{The Goal-Driven / Data-Driven binding, formally}

The architecture of the paper is now fully formal.  Let
\[
\mathcal{G} \;=\; \mathcal{A}_{\textsf{research}}
            \,\Vert\, \mathcal{A}_{\textsf{developer}}
            \,\Vert\, \mathcal{A}_{\textsf{mlebench}}
            \,\Vert\, \mathcal{A}_{\textsf{trainer}}
\]
under the orchestration relation $R_\mathcal{G}$ that allows any
of these sub-agents to consume the prefix emitted by any of the
others.  Then $\mathcal{G}$ is modelled by an idealised recursive-enumeration envelope.

Let
\[
\mathcal{V} \;=\; \mathcal{A}_{\textsf{tools\_index}}
            \,\Vert\, \mathcal{A}_{\textsf{unit\_test\_runner}}
\]
under the orchestration relation $R_\mathcal{V}$ that requires
every \textsf{test\_*} emission to be matched against the
\textsf{tools\_index} catalogue.  Then $L(\mathcal{V}) \in
\mathcal{L}_3$ for the verifier specification.

The validated trace language is
\[
L(\mathcal{G}\,\Vert\,\mathcal{V})
\;=\; L(\mathcal{G}) \,\cap\, L(\mathcal{V})
\]
which, by Theorem~\ref{thm:closure}, remains in the generator's model class, while conformance to $L(\mathcal{V})$ is decidable
in $O(n)$ on the verifier side.  The projected Type-2-style
sub-agents \texttt{planner\_agent}, \texttt{reviewer\_agent} and
the indexed-zone sub-agent \texttt{sop\_orchestrator} and the
Type-1-style sub-agents \texttt{task\_orchestrator},
\texttt{kernel\_pipeline} are not part of $\mathcal{G}$ or
$\mathcal{V}$ as designed; they are the \emph{projections} of
$\mathcal{G}\,\Vert\,\mathcal{V}$ onto the orchestration graph,
and they appear because the binding forces them into existence ---
not because the architect designed them.

The minimality claim is therefore:
\begin{quote}
\emph{One idealised goal-driven generator and one concrete Type-3
Data-Driven verifier suffice for the trace-validation architecture.
Everything else is a projected modelling view.}
\end{quote}

This is the operational meaning of ``a single Goal-Driven agent
tamed by the Data-Driven verifier results in success.''
